\section{The PoML Protocol}
\label{sec:protocol}

This section formally specifies the Proof-of-ML-Inference (PoML) consensus protocol. 

\subsection{System Model and Notation}
\label{sec:poml-system-def}

Let $\mathcal{N}$ denote the set of nodes (miners). Each node holds an identical model $M_\theta \in \mathcal{K}$ (Definition~\ref{def:model}) and a key pair $(\mathsf{pk}_m, \mathsf{sk}_m) \leftarrow \mathsf{KeyGen}_{\mathsf{pke}}(1^\kappa)$; we write $\mathcal{R}_{\mathsf{key}}$ for the induced key-pair relation, i.e., $(\mathsf{pk}_m, \mathsf{sk}_m) \in \mathcal{R}_{\mathsf{key}}$ iff $\mathsf{KeyGen}$ can output $(\mathsf{pk}_m, \mathsf{sk}_m)$. Since a public key may admit multiple valid secret keys, each node must also register a commitment to its secret key during registration:
\[
(c_{\mathsf{sk}},\, r_{\mathsf{sk}}) \leftarrow \mathsf{Commit}(\ppnew,\mathsf{sk}_m),
\]
and publish $c_{\mathsf{sk}}$ on-chain alongside~$\mathsf{pk}_m$. This binds the node to a unique secret key for the duration of its participation.

Let $\mathcal{U}$ denote the set of users submitting inference queries. Each user likewise holds a key pair $(\mathsf{pk}_u, \mathsf{sk}_u)$. A block is defined as $B = (s, x, \mathcal{X}, \mathcal{Y}, \Pi, \Pi^{(2)})$, where $s$ is the previous block hash, $x$ is the block content comprising standard coin transactions (transfers, fees, block rewards), $\mathcal{X}$ is the tuple of user queries to the model, $\mathcal{Y} = (\mathsf{ct}_1, \ldots, \mathsf{ct}_k)$ is the tuple of encrypted outputs, $\Pi = (\pi_1^{(1)}, \ldots, \pi_k^{(1)})$ and $\Pi^{(2)} = (\pi_1^{(2)}, \ldots, \pi_k^{(2)})$ ordered tuples of ZKPs which will be explained in the coming sections. 

We assume the following cryptographic primitives:
\begin{itemize}
    \item Commitment scheme: $(\mathsf{Setup}_\mathsf{com},\mathsf{Commit}, \open)$
  \item Hash function  $G,H : \{0,1\}^* \to \{0,1\}^\kappa$, used for seed derivation and the lottery
    \item Zero-knowledge proof systems $(\mathsf{Setup}_{\mathsf{zkp},1},\prover_1,\verifier_1)$, \\$(\mathsf{Setup}_{\mathsf{zkp},2},\prover_2,\verifier_2)$ for the two NP-languages described in the following sections
    \item Public-key encryption scheme $(\mathsf{KeyGen}_{\mathsf{pke}},\mathsf{Enc}, \mathsf{Dec})$
    \item Pseudorandom function $F : \{0,1\}^\kappa \times \{0,1\}^\kappa \to \{0,1\}^\kappa$, used to derive the prover's internal randomness 
    \item Digital signature scheme $(\mathsf{KeyGen}_{\mathsf{sig}}, \mathsf{Sign}, \mathsf{Verify}_{\mathsf{sig}})$
\end{itemize}

The public parameters $\mathsf{pp}$ and common reference string for the zero-knowledge proof $\mathsf{CRS}$ are assumed to be agreed out-of-band and published on the blockchain as parameters of the blockchain. The setup process for the commitment scheme and ZKP is also assumed to be done beforehand.

\subsection{Query Submission}

A user $u \in \mathcal{U}$ submits an inference request by providing a conditioning embedding $c = e(x)$, obtained by applying the public embedding function~$e$ to their input~$x$.

Each user maintains a running counter $\mathsf{taskID}$, which denotes the number of successful queries made. The user computes a commitment:
\begin{align}
    (\mathsf{c}_c, r_c) &\leftarrow \mathsf{Commit}(c \parallel \mathsf{taskID})
\end{align}

The user then signs and diffuses a query transaction to the network:
\[
\sigma_u \leftarrow \mathsf{Sign}(\mathsf{sk}_u,\, \mathsf{taskID} \parallel \mathsf{c}_c \parallel \text{fee}),
\]
\[
T_{\text{query}} = (\mathsf{pk}_u,\, \mathsf{taskID},\, \mathsf{c}_c,\, c,\, r_c,\, \text{fee},\, \sigma_u),
\]
where $\text{fee}$ is the cost of the query paid to the successful miner.
% The full transaction---including the embedding~$c$ and opening randomness~$r_c$---is diffused to the network so that miners can verify the commitment and perform inference. However, only the compact on-chain record
% \[
% \hat{T}_{\text{query}} = (\mathsf{pk}_u,\, \mathsf{taskID},\, \mathsf{c}_c,\, \text{fee},\, \sigma_u)
% \]
% is stored in the block. The embedding~$c$ and randomness~$r_c$ do not persist on-chain; the commitment~$\mathsf{c}_c$ suffices for proof verification.

\subsection{Query Retrieval}

Each node mines a block by retrieving queries from the mempool.

For each query, the miner verifies the user's signature \\$\mathsf{Verify}_{\mathsf{sig}}(\mathsf{pk}_u, \mathsf{taskID} \parallel \mathsf{c}_c \parallel \text{fee},\, \sigma_u) = 1$, checks the commitment opening $\open(\mathsf{pp}, \mathsf{c}_c, c \| \mathsf{taskID}, r_c) = 1$, and verifies that the $\mathsf{taskID}$ is correctly in order. If all checks pass, the miner proceeds with model inference.

\subsection{Inference and Proof Generation}

To build the block, miners first prepares the block content $G(s,x)$, then compute inferences and generate proofs one by one, forming a \textit{proof chain} $\Pi = (\pi_1^{(1)}, \ldots, \pi_k^{(1)})$. For the $i$-th query in the chain, the miner first derives a random seed:
\begin{equation}
r_i = H(G(s,x) \parallel \mathsf{c}_{c,i} \parallel \mathsf{taskID}_i \parallel \mathsf{pk}_m \parallel i),
\end{equation}
where $s$ is the previous block hash and $\mathsf{pk}_m$ is the miner's public key.

The model inference is computed as:
\begin{equation}
y_i = M_\theta(c_i, r_i).
\end{equation}

\paragraph{Output encryption and proof chain binding.}
Define the \emph{chain binding} for position~$i$ as:
\[
\mathsf{bind}_i = 
\begin{cases}
G(s,x) & \text{if } i = 1, \\
H(\pi_{i-1}^{(1)}) & \text{if } i \geq 2.
\end{cases}
\]
\thomas{Michele--We have to design the chain binding this way, cannot just use a position binding like H(1), H(2), otherwise this will break the proof-to-lottery bijection in section 5.3 - An adversary that e.g. runs two proofs in parallel for position 1 ($\pi_A^1,\pi_B^1$), two proofs in parallel for position 2 ($\pi_C^2,\pi_D^2$) can admit six different lottery chances ($\pi_A^1$),($\pi_A^1,\pi_C^2$), ($\pi_A^1,\pi_D^2$), ... for just four proofs}

The output is encrypted together with the chain binding and task ID under the querying user's public key:
\begin{align}
\label{eq:ct}
\mathsf{ct}_i \leftarrow \mathsf{Enc}(\mathsf{pk}_u,\, y_i \parallel \mathsf{bind}_i \parallel \mathsf{taskID}).
\end{align}

\paragraph{Proof statement.}
Nodes generate a zero-knowledge proof $\pi_i^{(1)}$ for membership in the NP language:

\[
\mathcal{L}_{\text{PoML}} 
= \left\{\,
  \begin{array}{@{}l@{}}
    (\mathsf{pp},\mathsf{taskID},\mathsf{bind}_i,\mathsf{c}_c,c_\theta, r_i, \mathsf{ct}_i, \mathsf{pk}_u):\\[0.2em]
    \exists\,(c,r_c,\theta,r_\theta,y) \ \text{s.t.} \\
    M_\theta(c, r_i)=y, \\
    \mathsf{Open}(\mathsf{pp},\mathsf{c}_c,c \,\|\, \mathsf{taskID},r_c) = 1, \\
    \mathsf{Open}(\mathsf{pp},c_\theta,\theta, r_\theta) = 1, \\
    \mathsf{ct}_i = \mathsf{Enc}(\mathsf{pk}_u,\, y \,\|\, \mathsf{bind}_i \,\|\, \mathsf{taskID})
  \end{array}
\right\}.
\]
We write $\mathsf{stmt}_i^{(1)} = (\mathsf{pp},\mathsf{taskID}_i,\mathsf{bind}_i,\mathsf{c}_{c,i},c_\theta, r_i, \mathsf{ct}_i, \mathsf{pk}_{u,i})$ for the public statement and $\mathsf{wit}_i^{(1)} = (c_i,r_{c,i},\theta,r_\theta,y_i)$ for the witness.

\paragraph{Deterministic proof generation.}
To prevent proof grinding---i.e., an adversary re-running the prover with different internal randomness to obtain distinct proofs for the same inference---we require that the prover's internal randomness is derived deterministically from the seed and the miner's secret key. Concretely, a pseudorandom function keyed by the miner's secret key expands the seed into prover randomness:
\[
\rho_i = F(\mathsf{sk}_m,\, r_i),
\]
and the miner generates the inference proof as \\ $\pi_i^{(1)} = \mathsf{Prove}(\mathsf{CRS},\, \mathsf{stmt}_i^{(1)},\, \mathsf{wit}_i^{(1)}\,;\; \rho_i)$, where $\rho_i$ replaces the prover's internal random tape. This ensures that each (miner, seed) pair $(\mathsf{sk}_m, r_i)$ yields a unique proof $\pi_i^{(1)}$, eliminating any advantage from re-randomization.

\paragraph{Meta-proof of deterministic proving.}
To enforce that the miner indeed used the prescribed randomness $\rho_i$, the miner additionally generates a \emph{meta-proof} $\pi_i^{(2)}$ for the language:
\[
\mathcal{L}_{\text{det}} = \left\{\,
  \begin{array}{@{}l@{}}
    (\mathsf{pp},\, r_i,\, \mathsf{stmt}_i^{(1)},\pi_i^{(1)},\, \mathsf{pk}_m,\, c_{\mathsf{sk}}):\\[0.2em]
    \exists\, (\mathsf{sk}_m,\, r_{\mathsf{sk}},\, \mathsf{wit}_i^{(1)}) \ \text{s.t.} \\
    (\mathsf{pk}_m, \mathsf{sk}_m) \in \mathcal{R}_{\mathsf{key}}, \\
    \mathsf{Open}(\mathsf{pp},\, c_{\mathsf{sk}},\, \mathsf{sk}_m,\, r_{\mathsf{sk}}) = 1, \\
    \rho_i = F(\mathsf{sk}_m,\, r_i), \\
    \pi_i^{(1)} = \mathsf{Prove}(\mathsf{CRS},\, \mathsf{stmt}_i^{(1)},\, \mathsf{wit}_i^{(1)}\,;\; \rho_i)
  \end{array}
\right\}.
\]
We write $\mathsf{stmt}_i^{(2)} = (\mathsf{pp}, r_i, \mathsf{stmt}_i^{(1)}, \pi_i^{(1)}, \mathsf{pk}_m, c_{\mathsf{sk}})$ for the public statement and $\mathsf{wit}_i^{(2)} = (\mathsf{sk}_m, r_{\mathsf{sk}}, \mathsf{wit}_i^{(1)})$ for the witness.
The meta-proof attests that $\pi_i^{(1)}$ was computed using the deterministic randomness $\rho_i = F(\mathsf{sk}_m, r_i)$, where $\mathsf{sk}_m$ is the miner's secret key corresponding to the public key~$\mathsf{pk}_m$ and matching the registered commitment~$c_{\mathsf{sk}}$. The zero-knowledge property of $\pi_i^{(2)}$ ensures that $\mathsf{sk}_m$ remains hidden. Each block therefore includes, for the $i$-th query, the pair $(\pi_i^{(1)}, \pi_i^{(2)})$.

\subsection{Block Construction and Consensus}

Nodes race to produce a valid block. To define a valid block, we introduce a lottery model which each node can run as soon as it has completed a proof within the block.

\paragraph{Lottery model.}

After generating proofs $\pi_i^{(1)}$ ($i \geq 1$), a lottery is executed based on the evolving block hash. Let $\Pi_i = (\pi_1^{(1)}, \ldots, \pi_i^{(1)})$ denote the inference proof tuple after the $i$-th proof. Since each proof $\pi_i^{(1)}$ is deterministically derived from the seed~$r_i$ (enforced by $\pi_i^{(2)}$), the lottery outcome is uniquely determined by the inference. We check:
\begin{equation}
H_i = H(G(s,x),\, \Pi_i) < D,
\end{equation} By adjusting $D$, we can tune the desired probability of lottery winning, hence the difficulty of the chain and other desired properties such as average block time.

If the above condition holds at any point, the miner can announce the block by publishing $B = (s, x, \mathcal{X}, \mathcal{Y}, \Pi, \Pi^{(2)})$. Each completed inference-proof pair yields exactly one lottery evaluation, motivating miners to process more queries and generate more proofs to increase their chance of winning.

\paragraph{Block validity.}

\begin{definition}[Valid Block]
\label{def:valid-block}
A block $B = (s,x,\mathcal{X}, \mathcal{Y}, \Pi, \Pi^{(2)})$ is valid,  iff the following conditions all hold:
\begin{enumerate}[nosep]
    \item \textbf{Lottery won:} $H(G(s,x),\, \Pi) < D$;
    \item \textbf{Non-empty proof tuples:} $|\mathcal{X}|=|\mathcal{Y}|=|\Pi|=|\Pi^{(2)}| \ge 1$;
    \item \textbf{Inference proofs verify:} $\mathsf{Verify}(\mathsf{CRS}, \mathsf{stmt}_i^{(1)}, \pi_i^{(1)}) = 1$ for all~$i$;
    \item \textbf{Meta-proofs verify:} $\mathsf{Verify}(\mathsf{CRS}, \mathsf{stmt}_i^{(2)}, \pi_i^{(2)}) = 1$ for all~$i$;
    \item \textbf{Transactions valid:} all coin transactions in $x$ are valid.
\end{enumerate}
\end{definition}

\subsection{Output Delivery}
\label{subsec:output-delivery}

In PoML, inference outputs are delivered privately to the querying user via public-key encryption. For each completed inference $y_i$, the miner encrypts the output under the user's public key (eq.\ref{eq:ct})
and includes the ciphertext $\mathsf{ct}_i$ in the encrypted-output tuple~$\mathcal{Y}$ of the block. Upon receiving the published block, the user decrypts the output using their secret key: $y_i \parallel \mathsf{bind}_i \parallel \mathsf{taskID} = \mathsf{Dec}(\mathsf{sk}_u,\, \mathsf{ct}_i)$. 

\subsection{Consensus Algorithm}

Algorithm~\ref{alg:poml-main} summarises the PoML block production procedure.

The protocol follows the longest-chain rule: upon receiving a new valid block, each node verifies all five conditions in Definition~\ref{def:valid-block} and adopts the longest valid chain. Ties are broken by the first block received.

\begin{algorithm}[t]
\caption{PoML Block Production}
\label{alg:poml-main}
\begin{algorithmic}[1]
\Require $B_{\mathsf{prev}}$ (previous block), $D$ (difficulty threshold), $\mathsf{CRS}$, $\mathsf{pp}$, $(\mathsf{pk}_m, \mathsf{sk}_m)$, $c_{\mathsf{sk}}$, $M_\theta$, $c_\theta$, $F$, $G$, $H$ as defined in \S\ref{sec:poml-system-def}
\State Retrieve query transactions from mempool; assemble block content~$x$
\State Initialize $s \gets H(B_{\mathsf{prev}})$, $\Pi \gets ()$, $\Pi^{(2)} \gets ()$, $\mathcal{X} \gets ()$, $\mathcal{Y} \gets ()$, $i \gets 0$
\Repeat
    \State $i \gets i + 1$
    \State Select next query 
    \Statex \qquad $T_i := (\mathsf{pk}_{u,i}, \mathsf{taskID}_i, \mathsf{c}_{c,i}, c_i, r_{c,i}, \text{fee}_i, \sigma_{u,i})$ 
    \Statex \qquad from the mempool
    \State Verify signature:
    \Statex \qquad $\mathsf{Verify}_{\mathsf{sig}}(\mathsf{pk}_{u,i},\, \mathsf{taskID}_i \| \mathsf{c}_{c,i} \| \text{fee}_i,\, \sigma_{u,i}) = 1$
    \State Verify commitment opening:
    \Statex \qquad $\open(\mathsf{pp},\, \mathsf{c}_{c,i},\, c_i \| \mathsf{taskID}_i,\, r_{c,i}) = 1$
    \State Derive seed $r_i \gets H(G(s,x) \parallel \mathsf{c}_{c,i} \parallel \mathsf{taskID}_i \parallel \mathsf{pk}_m \parallel i)$
    \State Compute inference $y_i \gets M_\theta(c_i, r_i)$
    \State Encrypt output $\mathsf{ct}_i \gets \mathsf{Enc}(\mathsf{pk}_{u,i},\, y_i \parallel \mathsf{bind}_i \parallel \mathsf{taskID}_i)$
    \State Derive prover randomness $\rho_i \gets F(\mathsf{sk}_m, r_i)$
    \State Generate inference proof:
    \Statex \qquad $\pi_i^{(1)} \gets \mathsf{Prove}(\mathsf{CRS},\, \mathsf{stmt}_i^{(1)},\, \mathsf{wit}_i^{(1)}\,;\; \rho_i)$
    \State Generate meta-proof:
    \Statex \qquad $\pi_i^{(2)} \gets \mathsf{Prove}(\mathsf{CRS},\, \mathsf{stmt}_i^{(2)},\, \mathsf{wit}_i^{(2)})$
    \State $\Pi \gets \Pi \,\|\, \pi_i^{(1)}$; $\Pi^{(2)} \gets \Pi^{(2)} \,\|\, \pi_i^{(2)}$; 
    \Statex \qquad $\mathcal{X} \gets \mathcal{X} \,\|\, \hat{T}_i$; $\mathcal{Y} \gets \mathcal{Y} \,\|\, \mathsf{ct}_i$
    \Statex \qquad \textit{// $\hat{T}_i = (\mathsf{pk}_{u,i}, \mathsf{taskID}_i, \mathsf{c}_{c,i}, \text{fee}_i, \sigma_{u,i})$ is the compact on-chain record}
    \State Compute $H_i \gets H(G(s,x),\, \Pi)$
\Until{$H_i < D$}
\State \textbf{Publish} block $B = (s, x, \mathcal{X}, \mathcal{Y}, \Pi, \Pi^{(2)})$
\end{algorithmic}
\end{algorithm}